\NeedsTeXFormat{LaTeX2e}
\ProvidesClass{SykSUCourseWork}[2026/03/21 v1.3]

\LoadClass[a4paper, 14pt]{extarticle}

% ==============================
% КОДИРОВКИ И ЯЗЫКИ
% ==============================
\RequirePackage[T2A]{fontenc}
\RequirePackage[utf8]{inputenc}
\RequirePackage[english,russian]{babel}
\RequirePackage{mathtext}
\renewcommand{\bfdefault}{bx}

% ==============================
% ПОЛЯ СТРАНИЦЫ
% ==============================
\RequirePackage[
left=3cm,
right=1cm,
top=2cm,
bottom=2cm,
headheight=15pt,
headsep=10pt,
includehead
]{geometry}

% ==============================
% МЕЖСТРОЧНЫЙ ИНТЕРВАЛ И АБЗАЦЫ
% ==============================
\RequirePackage{setspace}
\setstretch{1.2}

% ОТСТУП ПЕРВОЙ СТРОКИ 1 СМ ДЛЯ ВСЕГО ТЕКСТА
\setlength{\parindent}{1cm}
\setlength{\parskip}{0pt}

% ==============================
% ВЫРАВНИВАНИЕ ПО ШИРИНЕ
% ==============================
\RequirePackage{ragged2e}
\justifying

% ==============================
% ЗАГОЛОВКИ РАЗДЕЛОВ (ОТСТУП ПЕРВОЙ СТРОКИ 1 СМ)
% ==============================
\RequirePackage{titlesec}
\RequirePackage{needspace}

% SECTION - ГЛАВА
\titleformat{\section}
{\normalfont\bfseries\fontsize{16pt}{20pt}\selectfont\justifying}
{\thesection}
{0pt}
{\hspace{1cm}}  % ОТСТУП ПЕРВОЙ СТРОКИ 1 СМ
\titlespacing{\section}{0pt}{12pt}{0pt}

% НОВАЯ ГЛАВА НА НОВОЙ СТРАНИЦЕ
\let\oldsection\section
\renewcommand{\section}{%
	\clearpage%
	\oldsection%
}

% SUBSECTION - ПОДРАЗДЕЛ
\titleformat{\subsection}
{\normalfont\bfseries\fontsize{14pt}{18pt}\selectfont\justifying}
{\thesubsection}
{0pt}
{\hspace{1cm}}  % ОТСТУП ПЕРВОЙ СТРОКИ 1 СМ
\titlespacing{\subsection}{0pt}{6pt}{0pt}

% ПРОВЕРКА МЕСТА ДЛЯ SUBSECTION
\let\oldsubsection\subsection
\renewcommand{\subsection}{%
	\needspace{3\baselineskip}%
	\oldsubsection%
}

% SUBSUBSECTION - ПОДПОДРАЗДЕЛ
\titleformat{\subsubsection}
{\normalfont\bfseries\fontsize{14pt}{18pt}\selectfont\justifying}
{\thesubsubsection}
{0pt}
{\hspace{1cm}}  % ОТСТУП ПЕРВОЙ СТРОКИ 1 СМ
\titlespacing{\subsubsection}{0pt}{4pt}{0pt}

% ПРОВЕРКА МЕСТА ДЛЯ SUBSUBSECTION
\let\oldsubsubsection\subsubsection
\renewcommand{\subsubsection}{%
	\needspace{2\baselineskip}%
	\oldsubsubsection%
}

% ==============================
% ВВЕДЕНИЕ, ЗАКЛЮЧЕНИЕ, БИБЛИОГРАФИЯ, ПРИЛОЖЕНИЯ
% ==============================
\newcommand{\introduction}{%
	\clearpage
	\section*{\hspace{1cm}Введение}
	\addcontentsline{toc}{section}{Введение}
}

\newcommand{\conclusion}{%
	\clearpage
	\section*{\hspace{1cm}Заключение}
	\addcontentsline{toc}{section}{Заключение}
}

\newcommand{\bibliolist}{%
	\clearpage
	\section*{\hspace{1cm}Библиографический список}
	\addcontentsline{toc}{section}{Библиографический список}
}

\newcommand{\appendixsection}[1]{%
	\clearpage
	\section*{\hspace{1cm}#1}
	\addcontentsline{toc}{section}{#1}
}

% ==============================
% НАСТРОЙКА КОЛОНТИТУЛОВ И НУМЕРАЦИИ
% ==============================
\RequirePackage{fancyhdr}
\pagestyle{fancy}
\fancyhf{}
\fancyhead[R]{\thepage}
\renewcommand{\headrulewidth}{0pt}
\renewcommand{\footrulewidth}{0pt}

% ТИТУЛЬНЫЙ ЛИСТ И СОДЕРЖАНИЕ БЕЗ НОМЕРОВ
\AtBeginDocument{%
	\pagestyle{empty}
}

% СОДЕРЖАНИЕ
\let\oldtableofcontents\tableofcontents
\renewcommand{\tableofcontents}{%
	\clearpage
	\pagestyle{empty}%
	\pagenumbering{gobble}%
	\oldtableofcontents%
	\clearpage
	\pagenumbering{arabic}%
	\setcounter{page}{1}%
	\pagestyle{fancy}%
}

% ==============================
% НАСТРОЙКА СОДЕРЖАНИЯ (ОТСТУП ПЕРВОЙ СТРОКИ 1 СМ)
% ==============================
\RequirePackage{tocloft}

% ЗАГОЛОВОК "СОДЕРЖАНИЕ" - ВЫРАВНИВАНИЕ ПО ШИРИНЕ
\renewcommand{\cfttoctitlefont}{\normalfont\bfseries\fontsize{16pt}{20pt}\selectfont\justifying}
\renewcommand{\cftaftertoctitle}{\hfill}
\setlength{\cftbeforetoctitleskip}{0pt}
\setlength{\cftaftertoctitleskip}{20pt}

% ОТСТУП ПЕРВОЙ СТРОКИ 1 СМ ДЛЯ ВСЕХ СТРОК СОДЕРЖАНИЯ
\setlength{\cftsecindent}{1cm}
\setlength{\cftsubsecindent}{1cm}
\setlength{\cftsubsubsecindent}{1cm}

% ШИРИНА ДЛЯ НОМЕРОВ
\setlength{\cftsecnumwidth}{1.8cm}
\setlength{\cftsubsecnumwidth}{1.8cm}
\setlength{\cftsubsubsecnumwidth}{1.8cm}

% ШРИФТЫ С ВЫРАВНИВАНИЕМ ПО ШИРИНЕ
\renewcommand{\cftsecfont}{\normalfont\fontsize{12pt}{14pt}\selectfont\justifying}
\renewcommand{\cftsubsecfont}{\normalfont\fontsize{12pt}{14pt}\selectfont\justifying}
\renewcommand{\cftsubsubsecfont}{\normalfont\fontsize{12pt}{14pt}\selectfont\justifying}
\renewcommand{\cftsecpagefont}{\normalfont\fontsize{12pt}{14pt}\selectfont}
\renewcommand{\cftsubsecpagefont}{\normalfont\fontsize{12pt}{14pt}\selectfont}

% ТОЧКИ-ЛИДЕРЫ
\renewcommand{\cftsecleader}{\cftdotfill{\cftdotsep}}
\renewcommand{\cftsubsecleader}{\cftdotfill{\cftdotsep}}

% ==============================
% ПОДПИСИ РИСУНКОВ И ТАБЛИЦ
% ==============================
\RequirePackage{caption}
\RequirePackage{subcaption}

\captionsetup[figure]{
	labelsep=space,
	justification=centering,
	position=bottom,
	aboveskip=6pt,
	belowskip=0pt
}
\renewcommand{\figurename}{Рисунок}

\captionsetup[table]{
	labelsep=space,
	justification=centering,
	position=top,
	aboveskip=6pt,
	belowskip=0pt
}
\renewcommand{\tablename}{Таблица}

% ==============================
% ЛИСТИНГИ КОДА
% ==============================
\RequirePackage{listings}
\RequirePackage{xcolor}

\definecolor{codegray}{gray}{0.9}
\definecolor{backcolour}{rgb}{0.95,0.95,0.95}

\lstset{
	backgroundcolor=\color{backcolour},
	basicstyle=\ttfamily\fontsize{10pt}{12pt}\selectfont,
	breaklines=true,
	numbers=left,
	numberstyle=\tiny\color{codegray},
	frame=single,
	captionpos=t,
	tabsize=2,
	xleftmargin=0.3cm,
	xrightmargin=0.3cm
}
\renewcommand{\lstlistingname}{Листинг}

% ==============================
% МАТЕМАТИЧЕСКИЕ ФОРМУЛЫ
% ==============================
\RequirePackage{amsmath}
\RequirePackage{amssymb}

\AtBeginDocument{%
	\abovedisplayskip=12pt plus 3pt minus 3pt
	\abovedisplayshortskip=6pt plus 2pt minus 2pt
	\belowdisplayskip=12pt plus 3pt minus 3pt
	\belowdisplayshortskip=6pt plus 2pt minus 2pt
}

% ==============================
% ГРАФИКА, ТАБЛИЦЫ, ССЫЛКИ
% ==============================
\RequirePackage{graphicx}
\RequirePackage{float}
\RequirePackage{booktabs}
\RequirePackage{multirow}
\RequirePackage{array}
\RequirePackage{url}
\RequirePackage{hyperref}

\hypersetup{
	colorlinks=true,
	linkcolor=black,
	citecolor=black,
	urlcolor=blue
}

\endinput